\chapter{Evaluation}
\label{chapter:evaluation}

This chapter details the evaluation strategy that we are going to follow to assess the performance of ~\project{} and the results that we have obtained.
We start by presenting the research questions that we aim to answer with our evaluation. 
We then describe the datasets and benchmarks that we use for the evaluation and the experimental setup.
Following that, we present some results that we have obtained from the evaluation, and the plans to further evaluate the system.
Finally, we conclude the chapter by summarizing the key findings of the evaluation.

\section{Research Questions}
\label{eval:structural_research_questions}

So far we have presented the design and implementation with a special focus on making it competitive with the state-of-the-art graph analytics systems.
The objective has been to create a graph database that is as general as any standard graph database, but also as fast as any specialized graph analytics system.

In this chapter, we evaluate the performance of ~\project{} on structural graphs to answer the following research questions:
\begin{enumerate}
  \item How do the \textbf{EdgeKey and Adjacency List representations compare in terms of performance and memory usage}? We want to understand the trade-offs between the two representations and how they affect the performance of graph processing. We evaluate this using microbenchmarks that focus on the performance of basic graph operations such as edge insertion, edge deletion, and edge lookup.
  \item How much time do competitor systems spend in preprocessing the graphs? \textbf{Including preprocessing time, does~\project{} perform better?} What is a good metric to measure the "preprocessing time" in ~\project{} considering that we have no explicit preprocessing step?
  \item How does ~\project{} perform compared to other \textbf{in-memory graph processing systems} on standard graph analytics workloads? This corresponds to the case where the graph is small enough to fit in the \wt{} cache and the graph processing is done in memory.
  \item How does ~\project{} perform compared to other \textbf{out-of-core graph processing systems} on standard graph database workloads? This corresponds to the case where the graph is too large to fit in memory and some evictions from the \wt{} cache are necessary.
  \item How does ~\project{} perform compared to other \textbf{dynamic graph processing systems} on specialized graph analytics workloads? This corresponds to the case where the graph is updated frequently, and the graph processing is done in memory. We will assess the number of transactions that ~\project{} can handle per second. We will also evaluate the strong scaling of ~\project{} by increasing the number of threads and measuring the throughput.
  \item We need to assess the \textbf{overhead of storing graphs in our KV-store-compatible representations}. can we represent graphs in ~\project{} compared to other graph databases and out-of-core graph processing systems? If we enable compression, how does it affect the performance?
  \item How does the performance of ~\project{} change with \textbf{different graph sizes and types}? We know that the performance of graph processing systems is highly dependent on the graph structure. We want to understand how ~\project{} performs on different types of graphs (e.g., social networks, web graphs, uniform random graphs, and low-degree high-diameter graphs such as road networks).
  \item Can we use ~\project{}'s \api{} to write \textbf{edge-centric and vertex-centric algorithms}? Does one representation perform better than the other for these algorithms? We are not interested in measuring the performance of these algorithms, but rather in understanding how the \api{} can be used to implement them\footnote{Margo's dancing elephant analogy: "It's not about how well the bear dances, but that it dances at all."}.
\end{enumerate}

\section{Datasets and Benchmarks}
\label{eval:datasets_and_benchmarks}
A few graph analytics algorithms are commonly used in the graph processing community.
The most common of the graph algorithm "kernels" are formally specified and have a performant and \emph{correct} implementation provided by the GAP Benchmark Suite\ac{GAPBS}~\cite{DBLP:journals/corr/BeamerAP15}.
More importantly, though, using this reference implementation allows us to draw a more direct comparison between the performance of our graph database and other graph processing systems.
When running fair comparative benchmarks, it is important to execute the same algorithm implementation on all systems.
Every graph processing system has its own native implementation of graph kernels. 
Comparing these different implementations can lead to misleading conclusions because it is challenging to determine whether the observed performance is due to the inherent efficiency of the system or due to the use of a more efficient kernel implementation.

We use PageRank(PR), Breadth-First Search(BFS), Connected Components(CC), Single-Source Shortest Path (SSSP), and Triangle Counting(TC) algorithms from \ac{GAPBS} to evaluate the performance of ~\project{}.
These algorithms are widely used in the graph processing community and are representative of the different types of graph algorithms that are used in practice.

\subsection{Modifications to GAPBS}
\PM{this can be moved to the apppendix}
We use the \ac{GAPBS}~\cite{beamer2015gap} to run the graph algorithms on our graph database.
The GAP benchmark suite implementation has been widely used to compare graph processing systems, and therefore it is a performant and sensible choice for a more widespread comparison~\cite{teseo2021deLeo}.
While the graph API used in \ac{GAPBS} is standard and available in~\project{}, there are a few modifications we had to make to get \ac{GAPBS} working, which we describe in Appendix~\autoref{appendix:pagerank_graphapi}.

Another point to recapitulate here is that \project{} \emph{never} relabels nodes or edges.
We always use user-provided node IDs and edge IDs to store and access the graph.
This is a design decision that we made to simplify the graph representation and to avoid the overhead of maintaining a mapping between the user-provided IDs and the internal IDs.
This design decision has only a slight implication for our use of \ac{GAPBS} implementation, which assumes that the nodes are numbered from 0 to $MAX\_NODE\_ID$.
Instead of allocating our \code{pvector} and \code{sliding_queue} classes with a size equal to number of nodes, we instead initialize the classes with a buffer large enough to hold $MAX\_NODE\_ID$ number of nodes in the graph.
This is a minor change and does not affect the performance of the \ac{GAPBS} implementation, especially considering our use of \unix{mmap} regions to store the state.

\subsection{Datasets}
\label{eval:datasets}

\begin{table}[]
  \centering
  \begin{tabular}{|l|l|l|l|}
  \hline
  \textbf{Graphs} & \textbf{Vertices $|V|$} & \textbf{Edges $|E|$} & \textbf{\begin{tabular}[c]{@{}l@{}}Uncompressed\\ Size\end{tabular}} \\ \hline
  cit-Patents~\cite{citPatents}                                      & 3.77M   & 16.5M  &  268 MB \\ \hline
  dimacs9-USA~\cite{FullUSA92:online}                                & 23.94M  & 57.7M  & 940 MB \\ \hline
  twitter\_mpi~\cite{icwsm10cha}                                     & 52.57M  & 1.96B  &  37 GB\\ \hline
  \begin{tabular}[c]{@{}l@{}}kronecker-s26,\\ uniform-rand-s26\end{tabular} & 67.10M  & 1.07B  & 18 GB \\ \hline
  \begin{tabular}[c]{@{}l@{}}kronecker-s28,\\ uniform-rand-s28\end{tabular} & 268.43M & 4.29B  & 74 GB  \\ \hline
  \begin{tabular}[c]{@{}l@{}}kronecker-s30,\\ uniform-rand-s30\end{tabular} & 1.07B   & 17.18B &  \\ \hline
  \end{tabular}%
\caption{The graphs used for the evaluation with their number of vertices and edges, and the size of the uncompressed dataset.}
\label{tab:datasets}
\end{table}

We use the datasets mentioned in~\autoref{tab:datasets} for our evaluation.
The selection of the dataset is driven by the need to evaluate the performance of ~\project{} on different types of graphs.
We use the \textbf{dimacs9-USA} dataset~\cite{FullUSA92:online}, which is a road network graph of the USA. It is interesting because road network graphs are low-degree high-diameter graphs, and are known to be challenging for graph processing systems.
We use the \textbf{twitter\_mpi} dataset~\cite{icwsm10cha}, which is a social network graph. Social network graphs are known to have a power-law degree distribution and are predominantly used in graph processing system benchmarks.
We use the \textbf{kronecker} and \textbf{uniform-rand} datasets, which are synthetic graphs generated using the Kronecker graph generator.
These graphs are used to evaluate the performance of graph processing systems on graphs with controllably increasing sizes: from scale factor 26 (67.10M vertices, 1.07B edges) to scale factor 30 (1.07B vertices, 17.18B edges).
We use the Smooth Kronecker graph generator~\cite{anand2020smooth} to generate graphs with a power-law degree distribution, while the \ac{GAPBS} uniform-rand graph generator is used to generate graphs with a uniform degree distribution.

\section{Experimental Setup}
\label{eval:experimental_setup}
We conducted all machines on a single socket workstation equipped with an Intel(R) Xeon(R) W-2275 CPU @ 3.30GHz processor.
The workstation has 14 cores and 28 threads, 128GB of RAM, and a 1TB NVMe SSD.
We do not disable hyperthreading on the machine, and we use the default settings for the Linux kernel.
The source code was compiled with GCC 11.4.0 (which ships default on Ubuntu 22.04 LTS release) with the \code{-O3} optimization flag.
We use the WiredTiger v11.0.0 for the \project{} but set the \wt{} cache size to 10GB. This cache size is good enough for small graphs, but not large enough that the largest graph can fit in memory.

\section{Systems for Comparison}
\label{eval:systems_for_comparison}

\project{} is a graph database that can be used in specialized contexts for structural graph analytics, dynamic structural graph processing, and (most broadly) graph database workloads.
We span the whole grid of graph data management systems shown in Figure~\autoref{fig:distilled_design_space} and must therefor compare ~\project{} against all classes of graph processing systems.

In particular, we compare ~\project{} against the systems mentioned in~\autoref{tab:competitors}.
The Grafaree project mentioned in~\autoref{tab:competitors} is an in-house dynamic graph processing system benchmark framework that has been designed by Milad Rezaei, a former M.Sc. student at UBC.
It uses a custom read-write workload generator to evaluate the performance of dynamic graph processing systems under varying degrees of randomness of access.
Like Teseo, Grafaree also uses \ac{GAPBS} kernels to evaluate the analytics performance of dynamic graph processing systems.
We plan to use Grafaree to evaluate the performance of ~\project{} on dynamic graph processing workloads, and there is still some work to be done to get the \project{} driver ready for Grafaree (the timeline for this is discussed in ~\autoref{sec:future-timeline}).

\begin{table}[]
  \centering
  \resizebox{\textwidth}{!}{%
  \begin{tabular}{|l|l|l|l|l|l|}
  \hline
  \textbf{System Type} &
    \textbf{Benchmarks} &
    \textbf{How?} &
    \textbf{\begin{tabular}[c]{@{}l@{}}Relevant \project{}\\ Code Ready?\end{tabular}} &
    \textbf{Comparison Done?} &
    \textbf{Competitors} \\ \hline
  \textbf{OOC Analytics} &
    GAPBS kernels &
    GAPBS &
    {\color[HTML]{009901} \textbf{DONE}} &
    {\color[HTML]{CB0000} \textbf{NO}} &
    \begin{tabular}[c]{@{}l@{}}GraphChi~\cite{kyrola2012graphchi}, X-Stream~\cite{roy2013x}, \\ Lumos~\cite{vora2019lumos}, Blaze~\cite{kim2022blaze}\end{tabular} \\ \hline
  \textbf{In-Memory Analytics} &
    GAPBS kernels &
    GAPBS &
    {\color[HTML]{009901} \textbf{DONE}} &
    {\color[HTML]{F8A102} \textbf{IN-PROGRESS}} &
    GAPBS~\cite{beamer2015gap}, Ligra~\cite{shun2013ligra} \\ \hline
    \textbf{Dynamic graph stores} &
    \begin{tabular}[c]{@{}l@{}}Analytic Benchmarks, \\ Transactional workloads\end{tabular} &
    GAPBS + Grafaree &
    {\color[HTML]{F56B00} \textbf{IN-PROGRESS}} &
    {\color[HTML]{CB0000} \textbf{NO}} &
    \begin{tabular}[c]{@{}l@{}}GraphOne, Teseo,\\ LiveGraph\end{tabular} \\ \hline
  \textbf{Graph Databases} &
    \begin{tabular}[c]{@{}l@{}}LDBC queries,\\ Analytic queries\end{tabular} &
    GraphBenchmark~\cite{GraphDat14:online} &
    {\color[HTML]{CB0000} \textbf{NO}} &
    {\color[HTML]{CB0000} \textbf{NO}} &
    \begin{tabular}[c]{@{}l@{}}Neo4j~\cite{neo4j}/Arango~\cite{arangodb} \\ (single-node)\end{tabular} \\ \hline
  \end{tabular}%
  }
  \caption{Systems we compare against to evaluate the performance of ~\project{} on \textbf{structural graphs}. The systems are divided into four categories: Out-of-core(OOC) Analytics, In-memory Analytics, Dynamic graph stores, and Graph Databases.}
  \label{tab:competitors}
  \end{table}


\section{Microbenchmarks}
\label{eval:microbenchmarks}

We plan to use microbenchmarks to evaluate the performance of the EdgeKey and Adjacency List representations against each other and the overheads they introduce compared to a standard CSR representation.

The microbenchmarks will cover the following operations:
\begin{enumerate}
  \item \textbf{Edge Insertion/Deletion}: We will measure the time taken to insert or delete a single edge in the graph. We will also measure the time taken to insert/delete a batch of edges into the graph. This microbenchmark leverages the Grafaree workload generator to generate random edge insertions and deletions.
  \item \textbf{Vertex Insertion/Deletion}: We will measure the time taken to insert or delete a single vertex in the graph. We will also measure the time taken to insert/delete a batch of vertices into the graph. This microbenchmark leverages the Grafaree workload generator to generate random vertex insertions and deletions.
  \item \textbf{Vertex Seek}: This measures the time to find the vertex in the graph representation. We will measure the time taken to find a vertex by its ID. We will compare this to the time taken to find a vertex by its index in the CSR representation.
  \item \textbf{Neighborhood Scan}: This microbenchmark evaluates the time taken to scan the neighborhood of a randomly chosen vertex in the graph. We measure the time taken to scan the neighborhood and normalize it by the number of edges in the neighborhood, getting the average time per edge. We will compare this to the time taken to scan the neighborhood in the CSR representation.
\end{enumerate}

The microbenchmarks should give us a good idea of the overhead of using GraphAPI for accessing an in-memory graph compared to a CSR, which allows us to determine the overhead introduced by the use of GraphAPI and \wt{}. 

% \section{Graph Analytics Workloads}
% \label{eval:graph_analytics_workloads}

% \section{Graph Database Workloads}
% \label{eval:graph_database_workloads}

% \section{Dynamic Graph Processing Workloads}
% \label{eval:dynamic_graph_processing_workloads}


% \section{Storage Overhead}
% \label{eval:storage_overhead}

\section{Preliminary Results}
\label{eval:preliminary_results}
So far, the evaluation has focussed on ensuring that the \project{} project is competitive with the state-of-the-art in-memory graph processing benchmark: \ac{GAPBS}.
The idea here is to ensure that using~\project{} as the graph storage layer does not introduce a significant slowdown in the performance of the \emph{same} benchmark algorithm when run on an in-memory CSR representation.
Another important aspect of the evaluation is to demonstrate competitive performance with other out-of-core graph processing systems.
We have compared the performance of~\project{} with GraphChi~\cite{kyrola2012graphchi} on a few benchmarks that we report below.

\begin{table}[h!]
  \centering
  \resizebox{\textwidth}{!}{%
  \begin{tabular}{|l|l|l|l|}
  \hline
   & \textbf{\project{}}     & \textbf{\ac{GAPBS}}          & \textbf{GraphChi}      \\ \hline
  \multirow{2}{*}{\textbf{PageRank}}                  & \multirow{2}{*}{0.481 s} & 0.063 s (algo)      & 2.7s s (algo) \\ \cline{3-4} 
   &                         & 2.250 s (preprocessing) & 1.46 s (preprocessing) \\ \hline
  \multirow{2}{*}{\textbf{Triangle Counting (exact)}} & 3.77 s (Trust Triangle)  & \multirow{2}{*}{NA} & 2.448s (algo) \\ \cline{2-2} \cline{4-4} 
   & 0.067s (Cycle Triangle) &                         & 6.172s (preprocessing) \\ \hline
  \end{tabular}%
  }
  \caption{Preliminary benchmark performance data using the PageRank and Triangle Counting algorithms on the cit-Patents graph.}
  \label{tab:eval:structural_inmem}
  \end{table}

Table \autoref{tab:eval:structural_inmem} shows the preliminary performance data that we have obtained so far.
Flexograph performs better than GraphChi on the PageRank benchmark but is slower than the \ac{GAPBS} implementation.
\ac{GAPBS} is a highly optimized in-memory graph processing system, and it is not surprising that it outperforms ~\project{}.
However, \ac{GAPBS} does not provide an exact triangle counting implementation, and we had to use our own Trust Triangle and Cycle Triangle implementations, which are comparable with the algorithms used in GraphChi (Cycle triangles).
We significantly outperform GraphChi on this benchmark, which is a good sign for the performance of ~\project{} on structural graph analytics workloads.
